\documentclass[paper=a4paper,fontsize=10pt,jafontsize=9pt,titlepage]{jlreq}
%=================================================================
% LuaLaTeX preamble
% (%=================================================================
% LuaLaTeX preamble
% (%=================================================================
% LuaLaTeX preamble
% (\input{preamble}でtemplate.texに読み込まれることを想定)
%
% 用途:
%   日本語・英語・古典ギリシア語混在文書
%   (哲学・神学・古典学向け)
%
% 2026-06
%
% 重要:
%   - Unicode Greek を使用可能
%   - Babel/LGR の古い \textgreek{...} 入力を維持
%   - hyperref による PDF しおり対応
%=================================================================

%=================================================================
% 日本語処理
%=================================================================

\usepackage{luatexja-fontspec}
\usepackage{luatexja-ruby}

% 古典ギリシア語(Unicode)を和文扱いしない。
%
% これが無いと
%
%   λόγος
%   ἀλήθεια
%
% などが HaranoAjiMincho に送られ、
%
%   Missing character: There is no ό ...
%
% が発生する。
%
\ltjsetparameter{jacharrange={-2,-3}}

%=================================================================
% フォント
%=================================================================

% 欧文・ギリシア語
\setmainfont{Libertinus Serif}

% 和文
\setmainjfont{HaranoAjiMincho}

%=================================================================
% 古典ギリシア語
%=================================================================

% Babel の古い LGR 入力
%
%   \textgreek{l'ogos}
%   \textgreek{>'anjrwpoc}
%
% を維持するために必要。
%
% Unicode Greek
%
%   λόγος
%   \foreignlanguage{greek}{πάντες ἄνθρωποι}
%
% と共存可能。
%
%  しかしitalic, boldなどが出なくなるので、fontencは外す。その結果、LGRは不可
% \usepackage[LGR,T1]{fontenc}
\usepackage[polutonikogreek,english,japanese]{babel}

%=================================================================
% レイアウト
%=================================================================

\usepackage{paracol}
\footnotelayout{m}

\usepackage{longtable}

%=================================================================
% 色
%=================================================================

\usepackage{xcolor}

%=================================================================
% PDF ハイパーリンク
%=================================================================

\usepackage[
unicode,
bookmarks=true,
hidelinks
]{hyperref}

% hyperref の後に読むこと。
% 先に読むと Option clash が発生する。
%
\usepackage{bookmark}

%=================================================================
% ヘッダ・フッタ
%=================================================================

\usepackage{fancyhdr}
\pagestyle{fancy}
%=================================================================
% 注意事項
%=================================================================
%
% 見出し中のギリシア語:
%
%   \section{λόγος}
%
% は空白になる場合がある。
%
% その場合:
%
%   \section{{\rmfamily λόγος}}
%
% と書く。
%
%
% Babel ギリシア語:
%
%   \textgreek{l'ogos}
%
% Unicode ギリシア語:
%
%   \foreignlanguage{greek}{λόγος}
%
% の両方が利用可能。
%
%=================================================================
でtemplate.texに読み込まれることを想定)
%
% 用途:
%   日本語・英語・古典ギリシア語混在文書
%   (哲学・神学・古典学向け)
%
% 2026-06
%
% 重要:
%   - Unicode Greek を使用可能
%   - Babel/LGR の古い \textgreek{...} 入力を維持
%   - hyperref による PDF しおり対応
%=================================================================

%=================================================================
% 日本語処理
%=================================================================

\usepackage{luatexja-fontspec}
\usepackage{luatexja-ruby}

% 古典ギリシア語(Unicode)を和文扱いしない。
%
% これが無いと
%
%   λόγος
%   ἀλήθεια
%
% などが HaranoAjiMincho に送られ、
%
%   Missing character: There is no ό ...
%
% が発生する。
%
\ltjsetparameter{jacharrange={-2,-3}}

%=================================================================
% フォント
%=================================================================

% 欧文・ギリシア語
\setmainfont{Libertinus Serif}

% 和文
\setmainjfont{HaranoAjiMincho}

%=================================================================
% 古典ギリシア語
%=================================================================

% Babel の古い LGR 入力
%
%   \textgreek{l'ogos}
%   \textgreek{>'anjrwpoc}
%
% を維持するために必要。
%
% Unicode Greek
%
%   λόγος
%   \foreignlanguage{greek}{πάντες ἄνθρωποι}
%
% と共存可能。
%
%  しかしitalic, boldなどが出なくなるので、fontencは外す。その結果、LGRは不可
% \usepackage[LGR,T1]{fontenc}
\usepackage[polutonikogreek,english,japanese]{babel}

%=================================================================
% レイアウト
%=================================================================

\usepackage{paracol}
\footnotelayout{m}

\usepackage{longtable}

%=================================================================
% 色
%=================================================================

\usepackage{xcolor}

%=================================================================
% PDF ハイパーリンク
%=================================================================

\usepackage[
unicode,
bookmarks=true,
hidelinks
]{hyperref}

% hyperref の後に読むこと。
% 先に読むと Option clash が発生する。
%
\usepackage{bookmark}

%=================================================================
% ヘッダ・フッタ
%=================================================================

\usepackage{fancyhdr}
\pagestyle{fancy}
%=================================================================
% 注意事項
%=================================================================
%
% 見出し中のギリシア語:
%
%   \section{λόγος}
%
% は空白になる場合がある。
%
% その場合:
%
%   \section{{\rmfamily λόγος}}
%
% と書く。
%
%
% Babel ギリシア語:
%
%   \textgreek{l'ogos}
%
% Unicode ギリシア語:
%
%   \foreignlanguage{greek}{λόγος}
%
% の両方が利用可能。
%
%=================================================================
でtemplate.texに読み込まれることを想定)
%
% 用途:
%   日本語・英語・古典ギリシア語混在文書
%   (哲学・神学・古典学向け)
%
% 2026-06
%
% 重要:
%   - Unicode Greek を使用可能
%   - Babel/LGR の古い \textgreek{...} 入力を維持
%   - hyperref による PDF しおり対応
%=================================================================

%=================================================================
% 日本語処理
%=================================================================

\usepackage{luatexja-fontspec}
\usepackage{luatexja-ruby}

% 古典ギリシア語(Unicode)を和文扱いしない。
%
% これが無いと
%
%   λόγος
%   ἀλήθεια
%
% などが HaranoAjiMincho に送られ、
%
%   Missing character: There is no ό ...
%
% が発生する。
%
\ltjsetparameter{jacharrange={-2,-3}}

%=================================================================
% フォント
%=================================================================

% 欧文・ギリシア語
\setmainfont{Libertinus Serif}

% 和文
\setmainjfont{HaranoAjiMincho}

%=================================================================
% 古典ギリシア語
%=================================================================

% Babel の古い LGR 入力
%
%   \textgreek{l'ogos}
%   \textgreek{>'anjrwpoc}
%
% を維持するために必要。
%
% Unicode Greek
%
%   λόγος
%   \foreignlanguage{greek}{πάντες ἄνθρωποι}
%
% と共存可能。
%
%  しかしitalic, boldなどが出なくなるので、fontencは外す。その結果、LGRは不可
% \usepackage[LGR,T1]{fontenc}
\usepackage[polutonikogreek,english,japanese]{babel}

%=================================================================
% レイアウト
%=================================================================

\usepackage{paracol}
\footnotelayout{m}

\usepackage{longtable}

%=================================================================
% 色
%=================================================================

\usepackage{xcolor}

%=================================================================
% PDF ハイパーリンク
%=================================================================

\usepackage[
unicode,
bookmarks=true,
hidelinks
]{hyperref}

% hyperref の後に読むこと。
% 先に読むと Option clash が発生する。
%
\usepackage{bookmark}

%=================================================================
% ヘッダ・フッタ
%=================================================================

\usepackage{fancyhdr}
\pagestyle{fancy}
%=================================================================
% 注意事項
%=================================================================
%
% 見出し中のギリシア語:
%
%   \section{λόγος}
%
% は空白になる場合がある。
%
% その場合:
%
%   \section{{\rmfamily λόγος}}
%
% と書く。
%
%
% Babel ギリシア語:
%
%   \textgreek{l'ogos}
%
% Unicode ギリシア語:
%
%   \foreignlanguage{greek}{λόγος}
%
% の両方が利用可能。
%
%=================================================================

\lhead{{\itshape Summa Theologiae} I-II, q.67}
%--------------------

\bibliographystyle{jplain}

\title{{\bfseries PRIMA SECUNDAE}\\{\Huge Summae Theologiae}\\Sancti Thomae
Aquinatis\\{\sffamily QUEAESTIO SEXAGESIMASEPTIMA}\\DE DURATIONE VIRTUTUM POST HANC VITAM}
\author{Japanese translation\\by Yoshinori {\scshape Ueeda}}
\date{Last modified \today}

%%%% コピペ用
%\rhead{a.~}
%\begin{center}
% {\Large {\bfseries }}\\
% {\large }\\
% {\footnotesize }\\
% {\Large \\}
%\end{center}
%
%\begin{longtable}{p{21em}p{21em}}
%
%&
%
%
%\\
%
%\end{longtable}
%\newpage

\begin{document}

\maketitle
\thispagestyle{empty}

\begin{center}
{\LARGE 『神学大全』第二部の一}\\
{\Large 第六十七問\\この生の後の徳の持続について}
\end{center}

\begin{longtable}{p{21em}p{21em}}
{\scshape Deinde} considerandum est de duratione virtutum post hanc
 vitam. Et circa hoc quaeruntur sex.

 \begin{enumerate}
  \item utrum virtutes morales maneant post hanc vitam.
  \item utrum virtutes intellectuales.
  \item utrum fides.
  \item utrum remaneat spes.
  \item utrum aliquid fidei maneat, vel spei.
  \item utrum maneat caritas.
 \end{enumerate}

 &

 次にこの生の後の徳の持続について考察されるべきである。これを巡って六つのことが問われる。

 \begin{enumerate}
  \item 道徳的徳はこの生の後に留まるか。
  \item 知性的徳はどうか。
  \item 信仰はどうか。
  \item 希望は留まるか。
  \item 信仰の何かが留まるか、あるいは希望の何かがか。
  \item 愛徳は留まるか。
 \end{enumerate}
 
\end{longtable}


\newpage

\rhead{a.~1}
\begin{center}
{\Large {\bfseries ARTICULUS PRIMUS}}\\
{\large UTRUM VIRTUTES MORALES MANEANT POST HANC VITAM}\\
{\footnotesize II$^{a}$II$^{ae}$, q.134, a.1, ad 1; III {\itshape Sent.}, d.33, q.1, a.4; {\itshape De Virtut.}, q.5, a.4.}\\
{\Large 第一項\\道徳的徳はこの生の後に留まるか}
\end{center}

\begin{longtable}{p{21em}p{21em}}
 {\scshape Ad primum sic proceditur}. Videtur quod virtutes morales
 non maneant post hanc vitam. Homines enim in statu futurae gloriae
 erunt similes Angelis, ut dicitur Matth.~{\scshape xxii}. Sed
 ridiculum est in Angelis ponere virtutes morales, ut dicitur in X
 {\itshape Ethic}. Ergo neque in hominibus, post hanc vitam, erunt
 virtutes morales.
 
&

 第一項の問題へ、議論は以下のように進められる。道徳的徳はこの生の後は
 残らないと思われる。理由は以下の通り。人間は、『マタイによる福音書』
 第22章で言われるように、将来の栄光の状態においては天使に似たものとな
 る。しかるに『ニコマコス倫理学』第10巻で言われるように、天使の中に道
 徳的徳を措定するのは馬鹿げたことである。ゆえに人間の中にも、この生の
 後、道徳的徳はない。

\\





2.~{\scshape Praeterea}, virtutes morales perficiunt hominem in vita
activa. Sed vita activa non manet post hanc vitam, dicit enim
Gregorius, in VI {\itshape Moral}., {\itshape activae vitae opera cum
corpore transeunt}. Ergo virtutes morales non manent post hanc vitam.
 
&

 さらに、道徳的徳は活動的生において人間を完成する。しかるに活動的生が
 この生の後に留まることはない。というのもグレゴリウスが『道徳論』第4巻
 で「活動的生の業は身体と共に過ぎていく」と言っているからである。ゆえ
 に道徳的徳はこの生の後には留まらない。

\\





3.~{\scshape Praeterea}, temperantia et fortitudo, quae sunt virtutes
morales, sunt irrationalium partium, ut philosophus dicit, in III
{\itshape Ethic}. Sed irrationales partes animae corrumpuntur,
corrupto corpore, eo quod sunt actus organorum corporalium. Ergo
videtur quod virtutes morales non maneant post hanc vitam.
 
&

 さらに、節制と勇気は道徳的徳だが、哲学者が『ニコマコス倫理学』第3巻で
 言うように、非理性的部分に属する。しかるに魂の非理性的部分は身体が滅
 びると消滅する。なぜならそれらは身体的器官の現実態だからである。ゆえ
 に、道徳的徳はこの生の後には留まらない。

\\





 {\scshape Sed contra est} quod dicitur {\itshape Sap}.~{\scshape i},
 quod {\itshape iustitia perpetua est et immortalis}.
 
&

 しかし反対に、『知恵の書』第1章で「正義は永遠であり不死である」
 \footnote{「義は不死である。」(1:15)}と言われている。

\\

 {\scshape Respondeo dicendum} quod, sicut Augustinus dicit, in XIV
 {\itshape de Trin}., Tullius posuit post hanc vitam quatuor virtutes
 cardinales non esse; sed in alia vita homines {\itshape esse beatos
 sola cognitione naturae, in qua nihil est melius aut amabilius}, ut
 Augustinus dicit ibidem, ea {\itshape natura quae creavit omnes
 naturas}. Ipse autem postea determinat huiusmodi quatuor virtutes in
 futura vita existere, tamen alio modo.
 
&

 解答する。以下のように言われるべきである。アウグスティヌスが『三位一
 体論』第14巻で言うように、キケロは、この生の後に四つの枢要徳は存在し
 ないが、他の生において人間は、「そこにおいては何もそれ以上に善いもの
 や愛すべきものがない自然の認識だけによって幸福である」と考えた。その
 自然とは、アウグスティヌスが同じところで述べているように、「全ての自
 然を作った自然」である。他方、アウグスティヌス自身は、その後、そのよ
 うな四つの徳が将来の生においても存在するが、別の仕方によってであると
 結論している。

\\

 Ad cuius evidentiam, sciendum est quod in huiusmodi virtutibus
 aliquid est formale; et aliquid quasi materiale. Materiale quidem est
 in his virtutibus inclinatio quaedam partis appetitivae ad passiones
 vel operationes secundum modum aliquem. Sed quia iste modus
 determinatur a ratione, ideo formale in omnibus virtutibus est ipse
 ordo rationis.
 
&

 それを明らかにするために、以下のことが知られるべきである。このような
 徳においては形相的なものといわば質料的なものとがある。これらの徳にお
 いて質料的なものとは、欲求的部分の、ある限度に即した情念や働きへの傾
 向性である。しかしこの限度は理性によって限定されるので、全ての徳にお
 いて、理性の秩序それ自体は形相的なものである。

\\

 Sic igitur dicendum est quod huiusmodi virtutes morales in futura
 vita non manent, quantum ad id quod est materiale in eis. Non enim
 habebunt in futura vita locum concupiscentiae et delectationes
 ciborum et venereorum; neque etiam timores et audaciae circa pericula
 mortis; neque etiam distributiones et communicationes rerum quae
 veniunt in usum praesentis vitae.
 
&

 それゆえ、この意味で、このような道徳的徳は将来の生において、それらの
 中の質料的であるものに関しては残らないと言われるべきである。理由は以
 下の通り。将来の生には、食物やビーナス的なことへの欲情や愛好の場所が
 ないし、死の危険を巡る怖れや大胆さもない。また、現生の使用に役立つ諸
 事物の分配や伝達ということもない。

\\



 Sed quantum ad id quod est formale, remanebunt in beatis
 perfectissimae post hanc vitam, inquantum ratio uniuscuiusque
 rectissima erit circa ea quae ad ipsum pertinent secundum statum
 illum; et vis appetitiva omnino movebitur secundum ordinem rationis,
 in his quae ad statum illum pertinent.

 &

 しかし形相的なものに関しては、この生の後、最高に完全な徳が至福な人々
 の中に残るだろう。それは、各々のものの最高度に正しい理が、かの状態に
 おいてその人に属しているものを巡ってあり、また、かの状態に属するもの
 どもにおいて、欲求力があらゆる点で理性の秩序に従って動かされる限りに
 おいてである。

\\

 Unde Augustinus ibidem dicit quod {\itshape prudentia ibi erit sine
 ullo periculo erroris; fortitudo, sine molestia tolerandorum malorum;
 temperantia, sine repugnatione libidinum. Ut prudentiae sit nullum
 bonum Deo praeponere vel aequare; fortitudinis, ei firmissime
 cohaerere; temperantiae, nullo defectu noxio delectari.} De iustitia
 vero manifestius est quem actum ibi habebit, scilicet esse subditum
 Deo, quia etiam in hac vita ad iustitiam pertinet esse subditum
 superiori.
 
&

 このことからアウグスティヌスは同じ箇所で次のように述べている。「思慮
 はそこでどんな誤りの危険もなしにあるだろう。勇気は耐えられるべき悪の
 困難なしに、節制は情欲の反対なしにあるだろう。その結果、思慮にはどん
 な善も神に優先させたり神に並べたりすることがなく、勇気には神にこの上
 なく堅く結びつくことが属し、節制にはどんな有害な欠陥においても喜ぶこ
 とがない」。他方、正義については、そこでどんな作用を持つかがより明ら
 かである。すなわちそれは、神に従属するということであり、つまり、この
 生においてですら、上位のものに従属することは正義に属するからである。

\\

 {\scshape Ad primum ergo dicendum} quod philosophus loquitur ibi de
 huiusmodi virtutibus moralibus, quantum ad id quod materiale est in
 eis, sicut de iustitia, quantum ad {\itshape commutationes et
 depositiones}; de fortitudine, quantum ad {\itshape terribilia et
 pericula}; de temperantia, quantum ad {\itshape concupiscentias
 pravas}.
 
&

 第一異論に対しては、それゆえ、以下のように言われるべきである。哲学者
 はそこで、道徳的徳において質料的であるものに関して、それらについて語っ
 ている。すなわち正義については「交換と分配」に関して、勇気については
 「恐ろしいことと危険」について、節制については「歪められた欲情」に関
 してである。

\\


 Et similiter dicendum est ad secundum. Ea enim quae sunt activae
 vitae, materialiter se habent ad virtutes.

 
&

第二異論に対しても同様に言われるべきである。活動的な生に属するものは、
徳に対して質料的に関係する。

 
\\

 {\scshape Ad tertium dicendum} quod status post hanc vitam est
 duplex, unus quidem ante resurrectionem, quando animae erunt a
 corporibus separatae; alius autem post resurrectionem, quando animae
 iterato corporibus suis unientur.
 
&

 第三異論に対しては以下のように言われるべきである。この生の後の状態に
 は二通りある。一つは復活前であり、魂が身体から離れているときである。
 もう一つは復活後であり、魂が再び自分の身体に合一されたときである。

\\

 In illo ergo resurrectionis statu,
 erunt vires irrationales in organis corporis, sicut et nunc
 sunt. Unde et poterit in irascibili esse fortitudo, et in
 concupiscibili temperantia, inquantum utraque vis perfecte erit
 disposita ad obediendum rationi.
 
&

 ゆえに、かの復活の状態においては、今あるように、非理性的な力が身体の
 器官の中にあるであろう。したがって、怒情的部分には勇気が、欲情的部分
 には節制が、どちらの力も完全に理性に従うように態勢付けられている限り
 において、ありうるだろう。

\\

 Sed in statu ante resurrectionem,
 partes irrationales non erunt actu in anima, sed solum radicaliter in
 essentia ipsius, ut in primo dictum est. Unde nec huiusmodi virtutes
 erunt in actu nisi in radice, scilicet in ratione et voluntate, in
 quibus sunt seminalia quaedam harum virtutum, ut dictum est.

 
&

 しかし復活前の状態においては、非理性的部分は魂の中に現実態においては
 なく、第一部で言われたように、ただそれの本質の中に根のようなものとし
 てあるだろう。したがって、このような徳も、根においてでなければ、現実
 にはないであろう。根においてというのは、すでに述べられたとおり、理性
 と意志の中に、これらの徳の何らかの種子のようなものがあるという意味で
 ある。

\\


 Sed iustitia, quae est in voluntate, etiam actu remanebit. Unde
 specialiter de ea dictum est quod est perpetua et immortalis, tum
 ratione subiecti, quia voluntas incorruptibilis est; tum etiam
 propter similitudinem actus, ut prius dictum est.

 
&

 しかし正義は、意志の中にあるが、現実態においても残るだろう。このこと
 から、それについては特別に、永遠であり不死であると言われている。一つ
 にはそれは基体の観点からであり、というのも意志は不滅だから、もう一つ
 は働きの類似\footnote{上位のもの・神に従属すること。}のためであり、これは少し前に述べられた。

\end{longtable}
\newpage




\rhead{a.~2}
\begin{center}
{\Large {\bfseries ARTICULUS SECUNDUS}}\\
{\large UTRUM VIRTUTES INTELLECTUALES MANEANT POST HANCE VITAM}\\
{\footnotesize Part.~I, q.89, a.5, 6; III {\itshape Sent.}, d.31, q.2, a.4; IV, d.1, q.1, a.2; I {\itshape Cor.}, cap.13, lect.3.}\\
{\Large 第二項\\知性的徳はこの生の後も残るか}
\end{center}

\begin{longtable}{p{21em}p{21em}}

 {\scshape Ad secundum sic proceditur}. Videtur quod virtutes
 intellectuales non maneant post hanc vitam. Dicit enim apostolus, I
 {\itshape ad Cor}.~{\scshape xiii}, quod {\itshape scientia
 destruetur}, et ratio est quia {\itshape ex parte cognoscimus}. Sed
 sicut cognitio scientiae est ex parte, idest imperfecta; ita etiam
 cognitio aliarum virtutum intellectualium, quandiu haec vita
 durat. Ergo omnes virtutes intellectuales post hanc vitam cessabunt.

 
&

 第二項の問題へ、議論は以下のように進められる。知性的徳はこの生の後に
 は残らないと思われる。理由は以下の通り。使徒は『コリントの信徒への手
 紙一』で「学知は滅ぼされるだろう」\footnote{「愛は決して滅びません。
 しかし、預言は廃れ、異言はやみ、知識も廃れます。」(13:8)}と言うが、そ
 の理由は「私たちは部分的に思惟する」\footnote{「私たちの知識は一部分
 であり、預言も一部分だからです。」(13:9)}からである。しかるに学知の認
 識が部分的、すなわち不完全であるように、他の知性的徳の認識もまた、こ
 の生が続くかぎり不完全である。ゆえに全ての知性的徳はこの生の後は止む
 であろう。 

\\

2.~{\scshape Praeterea}, philosophus dicit, in {\itshape
Praedicamentis}, quod scientia, cum sit habitus, est qualitas
difficile mobilis, non enim de facili amittitur, nisi ex aliqua forti
transmutatione vel aegritudine. Sed nulla est tanta transmutatio
corporis humani sicut per mortem. Ergo scientia et aliae virtutes
intellectuales non manent post hanc vitam.


&

 さらに、哲学者は『カテゴリー論』で、学知は習慣なので動かされにくい性
 質であり、何らかの強い変容や病気によらないかぎり、容易には捨てられな
 いと言う。しかるに人間の身体にとって、死による変容ほど大きなものはな
 い。ゆえに学知と他の知性的徳は、この生の後には留まらない。

\\




3.~{\scshape Praeterea}, virtutes intellectuales perficiunt
intellectum ad bene operandum proprium actum. Sed actus intellectus
non videtur esse post hanc vitam, eo quod {\itshape nihil intelligit
anima sine phantasmate}, ut dicitur in III {\itshape de Anima};
phantasmata autem post hanc vitam non manent, cum non sint nisi in
organis corporeis. Ergo virtutes intellectuales non manent post hanc
vitam.
 
&

 さらに、知性的徳は、固有の作用を行うように知性を完成させる。しかるに
 知性の作用はこの生の後にはありえないと思われる。なぜなら、『デ・アニ
 マ』第3巻で言われるように「魂は表象像なしに何も認識しない」が、表象像
 は身体的器官の中にしかないから、この生の後には残らないからである。ゆ
 えに知性的な徳はこの生の後に残らない。

\\



 {\scshape Sed contra est} quod firmior est cognitio universalium et
 necessariorum, quam particularium et contingentium. Sed in homine
 remanet post hanc vitam cognitio particularium contingentium, puta
 eorum quae quis fecit vel passus est; secundum illud Luc.~{\scshape
 xvi}, {\itshape recordare quia recepisti bona in vita tua, et Lazarus
 similiter mala}. Ergo multo magis remanet cognitio universalium et
 necessariorum, quae pertinent ad scientiam et ad alias virtutes
 intellectuales.

 
&

 しかし反対に、普遍的で必然的な事柄にかんする認識は、個別的で非必然的
 な事柄の認識よりも堅固である。しかるに『ルカによる福音書』第16章「あ
 なたがあなたの人生の中でどんな善を受け取ったか思い出しなさい。そして
 ラザルスは同じように悪を受け取っていた」\footnote{「しかし、アブラハ
 ムは言った。『子よ、思い出すがよい。お前は生きている間に良いものを受
 け、ラザロのほうは悪いものを受けた。今は、ここで彼は慰められ、お前は
 もだえ苦しむのだ。」(16:25)}によれば、この生の後、人間の中に、人が行っ
 たことや被ったことなど、非必然的で個別的なことの認識が残る。ゆえにま
 してや、普遍的で必然的な事柄の認識は残るのであり、それらは学知や他の
 知性的徳に属する。

\\




 {\scshape Respondeo dicendum} quod, sicut in primo dictum est, quidam
 posuerunt quod species intelligibiles non permanent in intellectu
 possibili nisi quandiu actu intelligit, nec est aliqua conservatio
 specierum, cessante consideratione actuali, nisi in viribus
 sensitivis, quae sunt actus corporalium organorum, scilicet in
 imaginativa et memorativa. Huiusmodi autem vires corrumpuntur,
 corrupto corpore. Et ideo secundum hoc, scientia nullo modo post hanc
 vitam remanebit, corpore corrupto; neque aliqua alia intellectualis
 virtus.


&

 解答する。以下のように言われるべきである。第一巻で言われたように、あ
 る人々は、可知的形象は、現実態において知性認識をしている間しか、可能
 知性の中に留まらず、想像力や記憶力という身体的器官の現実態である感覚
 的なちからにおいてでなければ、現実的な考察が止むと形象の何らの保存も
 ないと論じた。しかるにこれらの力は身体が滅ぶと消滅する。ゆえにこの考
 えによると、学知はこの生の後に、身体が滅ぶと、どんな仕方によっても留
 まらず、それは他の知性的徳も同じである。

\\


 Sed haec opinio est contra sententiam Aristotelis, qui in III
 {\itshape de Anima} dicit quod {\itshape intellectus possibilis est
 in actu, cum fit singula, sicut sciens; cum tamen sit in potentia ad
 considerandum in actu}. 

&

 しかしこの意見はアリストテレスの論に反する。彼は『デ・アニマ』第3巻で、
 「可能知性は知るものとして個物となるときに現実態においてある。ただし、
 現実態において考察するためには可能態においてあるのだが」と述べている
 からである。

\\

Est etiam contra rationem, quia species intelligibiles recipiuntur in
 intellectu possibili immobiliter, secundum modum recipientis. Unde et
 intellectus possibilis dicitur {\itshape locus specierum}, quasi
 species intelligibiles conservans.

 &

 さらにこの意見は理性にも反する。というのも可知的形象は、受け取るもの
 のあり方に即して、不動の仕方で可能知性に受け取られるからである。した
 がって、可能知性は「形象の場所」と言われるが、それは可知的形象を保存
 するものという意味である。

 
 \\
 

 Sed phantasmata, ad quae respiciendo homo intelligit in hac vita,
 applicando ad ipsa species intelligibiles, ut in primo dictum est,
 corrupto corpore corrumpuntur. Unde quantum ad ipsa phantasmata, quae
 sunt quasi materialia in virtutibus intellectualibus, virtutes
 intellectuales destruuntur, destructo corpore, sed quantum ad species
 intelligibiles, quae sunt in intellectu possibili, virtutes
 intellectuales manent. Species autem se habent in virtutibus
 intellectualibus sicut formales. Unde intellectuales virtutes manent
 post hanc vitam, quantum ad id quod est formale in eis, non autem
 quantum ad id quod est materiale, sicut et de moralibus dictum est.

 
&

 しかし表象像は、この生において人間がそれへ目を向け、可知的形象をそれ
 へ適用することによって知性認識するのだが、第一部で述べられたように、
 身体が滅ぶと消滅する。したがって、知性的徳においていわば質料的なもの
 である表象像自体にかんしては、知性的徳は身体が滅ぶと破壊される。しか
 し可能知性の中にある可知的形象に関しては、知性的徳が残存する。しかる
 に形象は知性的徳において形相的なものとしてある。したがって知性的徳は
 この生の後に、それにおいて形相的であるものに関しては留まるが、質料的
 であるものに関しては留まらない。これは道徳的徳について述べられたこと
 と同様である。

\\




 {\scshape Ad primum ergo dicendum} quod verbum apostoli est
 intelligendum quantum ad id quod est materiale in scientia, et
 quantum ad modum intelligendi, quia scilicet neque phantasmata
 remanebunt, destructo corpore; neque erit usus scientiae per
 conversionem ad phantasmata.

 
&

 第一異論に対しては、それゆえ、以下のように言われるべきである。使徒の
 言葉は学知において質料的であるものに関して、そして知性認識のしかたに
 関して理解されるべきである。なぜなら、身体が壊れると表象像は残らない
 からであり、表象像に向き直ることによる学知の使用もないだろうからであ
 る。

\\




 {\scshape Ad secundum dicendum} quod per aegritudinem corrumpitur
 habitus scientiae quantum ad id quod est materiale in eo, scilicet
 quantum ad phantasmata, non autem quantum ad species intelligibiles,
 quae sunt in intellectu possibili.
 
&

 第二異論に対しては以下のように言われるべきである。病気によって学知の
 習慣はそれにおける質料的なもの、すなわち表象像に関して消滅する。しか
 し可能知性の中にある可知的形象に関しては消滅しない。

\\


 {\scshape Ad tertium dicendum} quod anima separata post mortem habet
 alium modum intelligendi quam per conversionem ad phantasmata, ut in
 primo dictum est. Et sic scientia manet, non tamen secundum eundem
 modum operandi, sicut et de virtutibus moralibus dictum est.

 
&

 第三異論に対しては以下のように言われるべきである。第一部で述べられた
 とおり、死後に離在する魂は表象像への向き直りによるのとは異なる知性認
 識のしかたを持つ。したがって学知は留まるが、道徳的徳について言われた
 のと同様、同じ働き方においてではない。
 
\\


\end{longtable}
\newpage



\rhead{a.~3}
\begin{center}
{\Large {\bfseries ARTICULUS TERTIUS}}\\
{\large UTRUM FIDES MANEAT POST HANC VITAM}\\
{\footnotesize II$^{a}$II$^{ae}$, q.4, a.4, ad 1; III {\itshape Sent.}, d.39, q.2, a.1, qu$^{a}$1; {\itshape De Virtut.}, q.5, a.4, ad 10.}\\
{\Large 第三項\\信仰はこの生の後に残るか}
\end{center}

\begin{longtable}{p{21em}p{21em}}



 {\scshape Ad tertium sic proceditur}. Videtur quod fides maneat post
 hanc vitam. Nobilior enim est fides quam scientia. Sed scientia manet
 post hanc vitam, ut dictum est. Ergo et fides.

 
&

 第三項の問題へ、議論は以下のように進められる。信仰はこの生の後にも残
 ると思われる。理由は以下の通り。信仰は学知よりも高貴である。しかるに、
 すでに述べられたとおり、学知はこの生の後も残る。ゆえに信仰も残る。

\\




 2.~{\scshape Praeterea}, I {\itshape ad Cor}.~{\scshape iii},
 dicitur, {\itshape fundamentum aliud nemo potest ponere, praeter id
 quod positum est, quod est Christus Iesus}, idest fides Christi
 Iesu. Sed sublato fundamento, non remanet id quod
 superaedificatur. Ergo, si fides non remanet post hanc vitam, nulla
 alia virtus remaneret.
 
&

 さらに、『コリントの信徒への手紙一』第1章で「イエス・キリストという置
 かれた基礎の他に、だれも他の基礎を据えることはできない」
 \footnote{「イエス・キリストというすでに据えられている土台のほかに、
 誰も他の土台を据えることはできないからです」(3:11)}と言われているが、
 この基礎はイエス・キリストへの信仰である。しかるに基礎が取り除かれる
 と、その上に建てられたものもなくなる。ゆえにもし信仰が残らないならば、
 他のどの徳も残らないだろう。

\\



3.~Praeterea, cognitio fidei et cognitio gloriae differunt secundum
perfectum et imperfectum. Sed cognitio imperfecta potest esse simul
cum cognitione perfecta, sicut in Angelo simul potest esse cognitio
vespertina cum cognitione matutina\footnote{Cf.~ST1, q.58, a.6, 7.};
et aliquis homo potest simul habere de eadem conclusione scientiam per
syllogismum demonstrativum, et opinionem per syllogismum
dialecticum\footnote{Cf.~{\itshape In I Post.~Analytic.}, l.31, n.4.:
Quia enim syllogismus dialecticus ad hoc tendit, ut opinionem faciat,
hoc solum est de intentione dialectici, ut procedat ex his, quae sunt
maxime opinabilia, et haec sunt ea, quae videntur vel pluribus, vel
maxime sapientibus. Et ideo si dialectico in syllogizando occurrat
aliqua propositio, quae secundum rei veritatem habeat medium, per quod
possit probari, sed tamen non videatur habere medium, sed propter sui
probabilitatem videatur esse per se nota; hoc sufficit dialectico, nec
inquirit aliud medium, licet propositio sit mediata, et, ex ea
syllogizans, sufficienter perficit dialecticum syllogismum. Sed
syllogismus demonstrativus ordinatur ad scientiam veritatis; et ideo
ad demonstratorem pertinet, ut procedat ex his, quae sunt secundum rei
veritatem immediata. Et si occurrat ei mediata propositio, necesse est
quod probet eam per medium proprium, quousque deveniat ad immediata,
nec est contentus probabilitate propositionis.}. Ergo etiam fides
simul esse potest, post hanc vitam, cum cognitione gloriae.
 
&

 さらに信仰の認識と栄光の認識は完全なものと不完全なものに即して異なる。
 しかるに不完全な認識は完全な認識と共にありうる。たとえば天使の中に、
 宵の認識と明けの認識が同時にありえ、
 また、ある人間が同一の結論について論証的三段論法による知識と、弁証的
 三段論法による意見とを同時に持ちうるように。ゆえに信仰もまた、この生
 の後に、栄光の認識と共にありうる。

\\




 {\scshape Sed contra est} quod apostolus dicit, {\itshape II ad
 Cor}.~{\scshape v}, {\itshape quandiu sumus in corpore, peregrinamur
 a domino, per fidem enim ambulamus, et non per speciem}. Sed illi qui
 sunt in gloria, non peregrinantur a domino, sed sunt ei
 praesentes. Ergo fides non manet post hanc vitam in gloria.
 
&

 しかし反対に、使徒は『コリントの信徒への手紙二』第5章で「私たちが身体
 の中にいる間、私たちは主から離れて歩いている。すなわち見ることによっ
 てではなく信仰によって私たちは歩いている」\footnote{「それで、私たち
 はいつも安心しています。もっとも、この体を住みかとしている間は、主か
 ら離れた身であることも知っています。／というのは、私たちは、直接見える
 姿によらず、信仰によって歩んでいるからです。」(5:6-7) }と言っている。
 しかし、栄光にある人々は主から離れて歩かず、主に現前している。ゆえに
 信仰はこの生の後、栄光においては留まらない。

 \\
 
 {\scshape Respondeo dicendum} quod oppositio est per se et propria
 causa quod unum excludatur ab alio, inquantum scilicet in omnibus
 oppositis includitur oppositio affirmationis et negationis.

 &

 解答する。以下のように言われるべきである。対立は、すべて対立するもの
 の中に肯定と否定の対立が含まれる限りにおいて、あるものが他のものか
 ら排除される自体的で固有の原因である。

 \\

 Invenitur autem in quibusdam oppositio secundum contrarias formas,
 sicut in coloribus album et nigrum.
 
&

 さて、あるものどもにおいて、対立は、反対する形相に即して見出される。
 たとえば、色における白と黒のように。

\\

 In quibusdam autem, secundum perfectum et imperfectum, unde in
 alterationibus magis et minus accipiuntur ut contraria, ut cum de
 minus calido fit magis calidum, ut dicitur in V {\itshape Physic}. Et
 quia perfectum et imperfectum opponuntur, impossibile est quod simul,
 secundum idem, sit perfectio et imperfectio.
 
&

 またあるものどもにおいては、完全と不完全に即して（対立が見出される）。
 このことから、『自然学』第5巻で言われるように、あまり熱くないものから
 より熱いものが生じる場合のように、変化において、より大きいものとより
 小さいものとが反対のものとして理解される。また、完全なものと不完全も
 のとは対立するので、同時に、同じ観点に即して、完全性と不完全性がある
 ことは不可能である。

\\


 Est autem considerandum quod imperfectio quidem quandoque est de
 ratione rei, et pertinet ad speciem ipsius, sicut defectus rationis
 pertinet ad rationem speciei equi vel bovis. Et quia unum et idem
 numero manens non potest transferri de una specie in aliam, inde est
 quod, tali imperfectione sublata, tollitur species rei, sicut iam non
 esset bos vel equus, si esset rationalis.
 
&

 しかし、たとえば理性を欠くことが馬や牛の種の性格に属する場合のように、
 不完全性は、事物の性格に属し、種それ自体に属することがある。そして一つ
 の数的に同一のものは、一つの種から別の種へと越えていくことはできない
 ので、そのような不完全性が取り除かれるならば、事物の種が破壊される。ちょ
 うど、もしそれらが理性的であったならば、それらはもはや牛でも馬でもな
 かったであろうように。

\\



 Quandoque vero imperfectio non pertinet ad rationem speciei, sed
 accidit individuo secundum aliquid aliud, sicut alicui homini
 quandoque accidit defectus rationis, inquantum impeditur in eo
 rationis usus, propter somnum vel ebrietatem vel aliquid
 huiusmodi. Patet autem quod, tali imperfectione remota, nihilominus
 substantia rei manet.
 
&

 他方、あるときには、不完全性が種の性格に属さず、むしろ何か或るものに
 即して個に附帯する。たとえば、ある人に、理性の欠陥が、睡眠や酩酊やそ
 のようなもののために、彼の中での理性の使用が妨げられる限りにおいて生
 起する場合のように。しかし、このような不完全性が取り除かれる場合、そ
 のような場合でも、事物の実体（本質）が留まることは明らかである。

\\

 Manifestum est autem quod imperfectio cognitionis est de ratione
 fidei. Ponitur enim in eius definitione, fides enim est {\itshape
 substantia sperandarum rerum, argumentum non apparentium}, ut dicitur
 ad {\itshape Heb}.~{\scshape xi}. Et Augustinus dicit, {\itshape Quid est fides?
 Credere quod non vides}. Quod autem cognitio sit sine apparitione vel
 visione, hoc ad imperfectionem cognitionis pertinet. Et sic
 imperfectio cognitionis est de ratione fidei. Unde manifestum est
 quod fides non potest esse perfecta cognitio, eadem numero manens.
 
&


 ところで、信仰の性格に認識の不完全性が属することは明らかである。なぜ
 ならその定義の中に、（不完全性が）置かれているからである。すなわち、
 『ヘブライ人への手紙』第11章で言われているように、信仰とは、「希望さ
 れるべき事物の実体、見えないものの証拠」である。また、アウグスティヌ
 スは「信仰とは何か。あなたが見ないものを信じることである」と述べてい
 る。しかるに、認識が現前や直視なしにあることは、認識の不完全性に属す
 る。そしてその意味で、認識の不完全性は信仰の性格に属する。したがって、
 信仰が、数において同一に留まりながら、完全な認識でありえないことは明
 らかである。
 

\\


 Sed ulterius considerandum est utrum simul possit esse cum cognitione
 perfecta, nihil enim prohibet aliquam cognitionem imperfectam simul
 esse aliquando cum cognitione perfecta. Est igitur considerandum quod
 cognitio potest esse imperfecta tripliciter, uno modo, ex parte
 obiecti cognoscibilis; alio modo, ex parte medii; tertio modo, ex
 parte subiecti.
 
&

 しかし、さらにそれが完全な認識と同時にありうるかが考察されるべきであ
 る。というのも、ある不完全な認識が、時として完全な認識と同時にあるこ
 とを妨げるものはないからである。ゆえに、認識は三つのしかたで不完全で
 ありうることが考察されるべきである。一つには、認識されうる対象の側か
 ら、もう一つには、媒介の側から、そして第三に、主体の側からである。

\\

 Ex parte quidem obiecti cognoscibilis, differunt secundum perfectum
 et imperfectum cognitio matutina et vespertina in Angelis, nam
 cognitio matutina est de rebus secundum quod habent esse in verbo;
 cognitio autem vespertina est de eis secundum quod habent esse in
 propria natura, quod est imperfectum respectu primi esse.

&

 認識されうる対象の側からは、天使における明けの認識と宵の認識が、完全
 なものと不完全なものに即して異なる。すなわち明けの認識は、言葉におい
 て存在を持つかぎりでの事物についてあるが、宵の認識は固有の本性におい
 て存在を持つかぎりにおけるそれらについてあるが、この存在は、第一の存
 在に比べると不完全である。

\\
 
 Ex parte vero medii, differunt secundum perfectum et imperfectum
 cognitio quae est de aliqua conclusione per medium demonstrativum, et
 per medium probabile.
 
&

 また媒介の側からは、ある結論について、論証的な媒介によってある認識と、
 蓋然的な媒介によってある認識が、完全なものと不完全なものに即して異な
 る。
 

\\
 

 Ex parte vero subiecti differunt secundum perfectum et imperfectum
 opinio, fides et scientia. Nam de ratione opinionis est quod
 accipiatur unum cum formidine alterius oppositi, unde non habet
 firmam inhaesionem. De ratione vero scientiae est quod habeat firmam
 inhaesionem cum visione intellectiva, habet enim certitudinem
 procedentem ex intellectu principiorum. Fides autem medio modo se
 habet, excedit enim opinionem, in hoc quod habet firmam inhaesionem;
 deficit vero a scientia, in hoc quod non habet visionem.
 
&

 また主体の側からは、意見、信仰、学知が、完全なものと不完全なものに即
 して異なる。すなわち、意見の性格には、ある事柄が、他方の反対者の怖れ
 を伴って受け取られていることが属するので、硬い結びつきがない。他方、
 学知の性格には、知性的な直視を伴う硬い結びつきをもつことが属する。な
 ぜなら、諸原理の知的認識から出てくる確実性を持つからである。そして信
 仰は、この中間のありかたをしている。つまり、硬い結びつきをもつことに
 おいて意見を超えているが、直視をもたないことにおいて学知より劣ってい
 るからである。
 

\\

 Manifestum est autem quod perfectum et imperfectum non possunt simul
 esse secundum idem, sed ea quae differunt secundum perfectum et
 imperfectum, secundum aliquid idem possunt simul esse in aliquo alio
 eodem.
 
 &

 ところで、同一のものに即して、完全なものと不完全なものは同時にありえ
 ないことは明らかだが、しかし、完全なものと不完全なものに即して異なる
 複数のものが、ある同一のことに即して、他のある同一のものの中に、同時
 にあることは可能である。

\\

 Sic igitur cognitio perfecta et imperfecta ex parte obiecti, nullo
 modo possunt esse de eodem obiecto. Possunt tamen convenire in eodem
 medio, et in eodem subiecto, nihil enim prohibet quod unus homo simul
 et semel per unum et idem medium habeat cognitionem de duobus, quorum
 unum est perfectum et aliud imperfectum, sicut de sanitate et
 aegritudine, et bono et malo.

 &

 ゆえにこのようにして、対象の側から見た完全な認識と不完全な認識は、けっ
 して同一の対象についてありえない。しかし、同一の媒介と同一の主体にお
 いて一致することはありうる。なぜなら、ある人が同時にある一つの時に、
 一つの同一の媒介によって、二つのものについての認識をもち、その一つは
 完全でもう一つは不完全であるということはありうるからである。たとえば、
 健康と病気についてや善と悪についてそうであるように。


\\

 Similiter etiam impossibile est quod cognitio perfecta et imperfecta
 ex parte medii, conveniant in uno medio. Sed nihil prohibet quin
 conveniant in uno obiecto, et in uno subiecto, potest enim unus homo
 cognoscere eandem conclusionem per medium probabile, et
 demonstrativum.
 
&

 さらに同様に、媒介の側から完全な認識と不完全な認識が、一つの媒介にお
 いて一致するというのも不可能である。しかし、一つの対象において、そし
 て一つの主体において、一致することを妨げるものはない。たとえば、ある
 人が、同一の結論を、蓋然的な媒介と論証的な媒介によって認識することは
 可能である。

\\

 Et est similiter impossibile quod cognitio perfecta et imperfecta ex
 parte subiecti, sint simul in eodem subiecto. Fides autem in sui
 ratione habet imperfectionem quae est ex parte subiecti, ut scilicet
 credens non videat id quod credit, beatitudo autem de sui ratione
 habet perfectionem ex parte subiecti, ut scilicet beatus videat id
 quo beatificatur, ut supra dictum est. Unde manifestum est quod
 impossibile est quod fides maneat simul cum beatitudine in eodem
 subiecto.

 &

 そして同様に、主体の側から完全な認識と不完全な認識が、同一の主体にお
 いて同時にあることは不可能である。しかし、信仰はその性格の中に、主体
 の側に由来する不完全性を持っている。すなわち、信じる者は、信じている
 事柄を見ていない。他方、至福は、その性格の中に主体の側に由来する完全
 性をもつ。なぜなら、至福者は、前に述べられたとおり、それによって至福
 とされているところのものを見ているからである。したがって、信仰が、同
 一の主体の中に、至福とともに留まることは不可能である。


\\

 {\scshape Ad primum ergo dicendum} quod fides est nobilior quam
 scientia, ex parte obiecti, quia eius obiectum est veritas prima. Sed
 scientia habet perfectiorem modum cognoscendi, qui non repugnat
 perfectioni beatitudinis, scilicet visioni, sicut repugnat ei modus
 fidei.

 
 &

 第一異論に対しては、それゆえ、以下のように言われるべきである。信仰は、
 対象という点から見ると、学知よりも高貴である。なぜなら、それの対象は
 第一真理であるから。しかし、学知はより完全な認識の様態をもっている。
 その様態は、信仰の様態がそれに反対するようには、直視という至福の完全
 性に反対しない。

\\

 {\scshape Ad secundum dicendum} quod fides est fundamentum quantum ad
 id quod habet de cognitione. Et ideo quando perficietur cognitio,
 erit perfectius fundamentum.
 
&

 第二異論に対しては以下のように言われるべきである。信仰は、認識につい
 てあるものにかんして、基礎である。ゆえに認識が完全にされるほど、基礎
 はより完全であるだろう。

\\


 {\scshape Ad tertium} patet solutio ex his quae dicta sunt.
 
 &

 第三異論に対しては、すでに述べられたことから解答が明らかである。



\end{longtable}
\newpage




\rhead{a.~4}
\begin{center}
{\Large {\bfseries ARTICULUS QUARTUS}}\\
{\large UTRUM SPES MANEAT POST MORTEM IN STATU GLORIAE}\\
{\footnotesize II$^{a}$II$^{ae}$, q.18, a.2; III {\itshape Sent.}, dist.26, q.2, a.5, qu$^{a}$2; dist.31, q.2, a.1, qu$^{a}$2; {\itshape De Virtut.}, q.4, a.4.}\\
{\Large 第四項\\希望は死後の栄光の状態において留まるか}
\end{center}

\begin{longtable}{p{21em}p{21em}}
 {\scshape Ad quartum sic proceditur}. Videtur quod spes maneat post
 mortem in statu gloriae. Spes enim nobiliori modo perficit appetitum
 humanum quam virtutes morales. Sed virtutes morales manent post hanc
 vitam, ut patet per Augustinum, in XIV {\itshape de Trin}. Ergo multo
 magis spes.

 &

 第四項の問題へ、議論は以下のように進められる。希望は、死後の栄光の状
 態においても留まると思われる。理由は以下の通り。希望は、道徳的諸徳よ
 りも高貴なしかたで、人間の欲求を完成する。しかるに道徳的諸徳は『三位
 一体論』第14巻のアウグスティヌスによって明らかなとおり、この生の後に
 も残る。ゆえに、ましてや、希望は残る。


\\

2. {\scshape Praeterea}, spei opponitur timor. Sed timor manet post
hanc vitam, et in beatis quidem timor filialis, qui manet in saeculum;
et in damnatis timor poenarum. Ergo spes, pari ratione, potest
permanere.
 

 &

 さらに、希望には怖れが対立する。しかるに怖れはこの生の後にも残り、至
 福者においては息子の怖れが世々に残り、悪霊においては罪の怖れが残る。
 ゆえに希望も同じ理由で留まりうる。


 \\
 
3. {\scshape Praeterea}, sicut spes est futuri boni, ita et
desiderium. Sed in beatis est desiderium futuri boni, et quantum ad
gloriam corporis, quam animae beatorum desiderant, ut dicit
Augustinus, XII {\itshape Super Gen.~ad Litt}.; et etiam quantum ad
gloriam animae, secundum illud {\itshape Eccli}.~{\scshape xxiv},
{\itshape qui edunt me, adhuc esurient, et qui bibunt me, adhuc
sitient}; et I Petr.~{\scshape i}, dicitur, {\itshape in quem
desiderant Angeli prospicere}. Ergo videtur quod possit esse spes post
hanc vitam in beatis.
 
 &

 さらに、希望が将来の善に属するように、願望もまたそうである。しかるに
 至福者の中には将来の善への願望がある。それは、アウグスティヌスが『創
 世記逐語注解』第12巻で言うように、身体の栄光に関してであり、さらには、
 『シラ書』第24章「私を食べる者はさらに飢え、私を飲む者はさらに渇く」
 \footnote{「私を食べる者は、さらに飢え／私を飲む者は、さらに渇く」。
 (29:12)}や、『ペトロの手紙一』第1章「天使たちはその人に見ることを欲し
 ている」\footnote{「天使たちも、うかがい見たいと願っていることなので
 す。」(1:12)}によれば、魂の栄光に関してでもある。ゆえにこの生の後、至
 福者の中にも希望はありうると思われる。

\\


 {\scshape Sed contra est} quod apostolus dicit, {\itshape
 Rom}.~{\scshape viii}, {\itshape quod videt quis, quid sperat?} Sed
 beati vident id quod est obiectum spei, scilicet Deum. Ergo non
 sperant.
 
 &

 しかし反対に、使徒は『ローマの信徒への手紙』第8章で「だれが見ているも
 のを希望するのか」\footnote{「現に見ているものを、誰がなお望むでしょ
 うか。」(8:24)}言っている。しかるに至福者たちは希望の対象であるもの、
 すなわち神を見ている。ゆえに彼らはそれを希望しない。

\\

 {\scshape Respondeo dicendum} quod, sicut dictum est, id quod de
 ratione sui importat imperfectionem subiecti, non potest simul stare
 cum subiecto opposita perfectione perfecto.

 &

 解答する。以下のように言われるべきである。すでに述べられたとおり、自
 らの性格の中に主体の不完全性を含意するものは、対立する完全性によって
 完全である主体と同時に立つことはできない。

 \\

 Sicut patet quod motus in
 ratione sui importat imperfectionem subiecti, est enim actus
 existentis in potentia, inquantum huiusmodi, unde quando illa
 potentia reducitur ad actum, iam cessat motus; non enim adhuc
 albatur, postquam iam aliquid factum est album.

 &

 たとえば、運動は、自らの性格の中に主体の不完全性を含意する。運動は、
 運動であるかぎり、可能態にあるものの現実態だからである。したがって、
 その可能態が現実態に引き出されるときには、すでに運動であることを止め
 る。たとえば、何かが白くされた後には、漂白の運動はすでにない。

 \\

 Spes autem importat motum quendam in id quod non habetur; ut patet ex
 his quae supra de passione spei diximus. Et ideo quando habebitur id
 quod speratur, scilicet divina fruitio, iam spes esse non poterit.
 
&

 しかるに希望は、私たちが前に希望の情念について述べたことから明らかな
 とおり、持たれていないものへの一種の運動を含意する。ゆえに希望された
 もの、すなわち神を享受することが持たれたとき、希望はもう存在すること
 ができないであろう。

\\


 {\scshape Ad primum ergo dicendum} quod spes est nobilior virtutibus
 moralibus quantum ad obiectum, quod est Deus. Sed actus virtutum
 moralium non repugnant perfectioni beatitudinis, sicut actus spei;
 nisi forte ratione materiae, secundum quam non manent. Non enim
 virtus moralis perficit appetitum solum in id quod nondum habetur;
 sed etiam circa id quod praesentialiter habetur.
 
&

 

\\





[36318] Iª-IIae q. 67 a. 4 ad 2
 Ad secundum dicendum quod timor est duplex, servilis et filialis, ut infra dicetur. Servilis quidem est timor poenae, qui non poterit esse in gloria, nulla possibilitate ad poenam remanente. Timor vero filialis habet duos actus, scilicet revereri Deum, et quantum ad hunc actum manet; et timere separationem ab ipso, et quantum ad hunc actum non manet. Separari enim a Deo habet rationem mali, nullum autem malum ibi timebitur, secundum illud Proverb. I, abundantia perfruetur, malorum timore sublato. Timor autem opponitur spei oppositione boni et mali, ut supra dictum est, et ideo timor qui remanet in gloria, non opponitur spei. In damnatis autem magis potest esse timor poenae, quam in beatis spes gloriae. Quia in damnatis erit successio poenarum, et sic remanet ibi ratio futuri, quod est obiectum timoris, sed gloria sanctorum est absque successione, secundum quandam aeternitatis participationem, in qua non est praeteritum et futurum, sed solum praesens. Et tamen nec etiam in damnatis est proprie timor. Nam sicut supra dictum est, timor nunquam est sine aliqua spe evasionis, quae omnino in damnatis non erit. Unde nec timor; nisi communiter loquendo, secundum quod quaelibet expectatio mali futuri dicitur timor.

 
&


\\





[36319] Iª-IIae q. 67 a. 4 ad 3
 Ad tertium dicendum quod quantum ad gloriam animae, non potest esse in beatis desiderium, secundum quod respicit futurum, ratione iam dicta. Dicitur autem ibi esse esuries et sitis, per remotionem fastidii, et eadem ratione dicitur esse desiderium in Angelis. Respectu autem gloriae corporis, in animabus sanctorum potest quidem esse desiderium, non tamen spes, proprie loquendo, neque secundum quod spes est virtus theologica, sic enim eius obiectum est Deus, non autem aliquod bonum creatum; neque secundum quod communiter sumitur. Quia obiectum spei est arduum, ut supra dictum est, bonum autem cuius iam inevitabilem causam habemus, non comparatur ad nos in ratione ardui. Unde non proprie dicitur aliquis qui habet argentum, sperare se habiturum aliquid quod statim in potestate eius est ut emat. Et similiter illi qui habent gloriam animae, non proprie dicuntur sperare gloriam corporis; sed solum desiderare.

 
&



 
\end{longtable}
\newpage


\end{document}



%\rhead{a.~}
%\begin{center}
% {\Large {\bfseries }}\\
% {\large }\\
% {\footnotesize }\\
% {\Large \\}
%\end{center}
%
%\begin{longtable}{p{21em}p{21em}}
%
%&
%
%
%\\
%
%\end{longtable}
%\newpage


ARTICULUS 5
[36320] Iª-IIae q. 67 a. 5 arg. 1
Ad quintum sic proceditur. Videtur quod aliquid fidei vel spei remaneat in gloria. Remoto enim eo quod est proprium, remanet id quod est commune, sicut dicitur in libro de causis, quod, remoto rationali, remanet vivum; et remoto vivo, remanet ens. Sed in fide est aliquid quod habet commune cum beatitudine, scilicet ipsa cognitio, aliquid autem quod est sibi proprium, scilicet aenigma; est enim fides cognitio aenigmatica. Ergo, remoto aenigmate fidei, adhuc remanet ipsa cognitio fidei.

[36321] Iª-IIae q. 67 a. 5 arg. 2
Praeterea, fides est quoddam spirituale lumen animae, secundum illud Ephes. I, illuminatos oculos cordis vestri in agnitionem Dei; sed hoc lumen est imperfectum respectu luminis gloriae, de quo dicitur in Psalmo XXXV, in lumine tuo videbimus lumen. Lumen autem imperfectum remanet, superveniente lumine perfecto, non enim candela extinguitur, claritate solis superveniente. Ergo videtur quod ipsum lumen fidei maneat cum lumine gloriae.

[36322] Iª-IIae q. 67 a. 5 arg. 3
Praeterea, substantia habitus non tollitur per hoc quod subtrahitur materia, potest enim homo habitum liberalitatis retinere, etiam amissa pecunia; sed actum habere non potest. Obiectum autem fidei est veritas prima non visa. Ergo, hoc remoto per hoc quod videtur veritas prima, adhuc potest remanere ipse habitus fidei.

[36323] Iª-IIae q. 67 a. 5 s. c.
Sed contra est quod fides est quidam habitus simplex. Simplex autem vel totum tollitur, vel totum manet. Cum igitur fides non totaliter maneat, sed evacuetur, ut dictum est; videtur quod totaliter tollatur.

[36324] Iª-IIae q. 67 a. 5 co.
Respondeo dicendum quod quidam dixerunt quod spes totaliter tollitur, fides autem partim tollitur, scilicet quantum ad aenigma, et partim manet, scilicet quantum ad substantiam cognitionis. Quod quidem si sic intelligatur quod maneat non idem numero, sed idem genere, verissime dictum est, fides enim cum visione patriae convenit in genere, quod est cognitio. Spes autem non convenit cum beatitudine in genere, comparatur enim spes ad beatitudinis fruitionem, sicut motus ad quietem in termino. Si autem intelligatur quod eadem numero cognitio quae est fidei, maneat in patria; hoc est omnino impossibile. Non enim, remota differentia alicuius speciei, remanet substantia generis eadem numero, sicut, remota differentia constitutiva albedinis, non remanet eadem substantia coloris numero, ut sic idem numero color sit quandoque albedo, quandoque vero nigredo. Non enim comparatur genus ad differentiam sicut materia ad formam, ut remaneat substantia generis eadem numero, differentia remota; sicut remanet eadem numero substantia materiae, remota forma. Genus enim et differentia non sunt partes speciei, alioquin non praedicarentur de specie. Sed sicut species significat totum, idest compositum ex materia et forma in rebus materialibus, ita differentia significat totum, et similiter genus, sed genus denominat totum ab eo quod est sicut materia; differentia vero ab eo quod est sicut forma; species vero ab utroque. Sicut in homine sensitiva natura materialiter se habet ad intellectivam, animal autem dicitur quod habet naturam sensitivam; rationale quod habet intellectivam; homo vero quod habet utrumque. Et sic idem totum significatur per haec tria, sed non ab eodem. Unde patet quod, cum differentia non sit nisi designativa generis, remota differentia, non potest substantia generis eadem remanere, non enim remanet eadem animalitas, si sit alia anima constituens animal. Unde non potest esse quod eadem numero cognitio, quae prius fuit aenigmatica, postea fiat visio aperta. Et sic patet quod nihil idem numero vel specie quod est in fide, remanet in patria; sed solum idem genere.

[36325] Iª-IIae q. 67 a. 5 ad 1
Ad primum ergo dicendum quod, remoto rationali, non remanet vivum idem numero, sed idem genere, ut ex dictis patet.

[36326] Iª-IIae q. 67 a. 5 ad 2
Ad secundum dicendum quod imperfectio luminis candelae non opponitur perfectioni solaris luminis, quia non respiciunt idem subiectum. Sed imperfectio fidei et perfectio gloriae opponuntur ad invicem, et respiciunt idem subiectum. Unde non possunt esse simul, sicut nec claritas aeris cum obscuritate eius.

[36327] Iª-IIae q. 67 a. 5 ad 3
Ad tertium dicendum quod ille qui amittit pecuniam, non amittit possibilitatem habendi pecuniam, et ideo convenienter remanet habitus liberalitatis. Sed in statu gloriae non solum actu tollitur obiectum fidei, quod est non visum; sed etiam secundum possibilitatem, propter beatitudinis stabilitatem. Et ideo frustra talis habitus remaneret.


%\rhead{a.~}
%\begin{center}
% {\Large {\bfseries }}\\
% {\large }\\
% {\footnotesize }\\
% {\Large \\}
%\end{center}
%
%\begin{longtable}{p{21em}p{21em}}
%
%&
%
%
%\\
%
%\end{longtable}
%\newpage


ARTICULUS 6
[36328] Iª-IIae q. 67 a. 6 arg. 1
Ad sextum sic proceditur. Videtur quod caritas non maneat post hanc vitam in gloria. Quia, ut dicitur I ad Cor. XIII, cum venerit quod perfectum est, evacuabitur quod ex parte est, idest quod est imperfectum. Sed caritas viae est imperfecta. Ergo evacuabitur, adveniente perfectione gloriae.

[36329] Iª-IIae q. 67 a. 6 arg. 2
Praeterea, habitus et actus distinguuntur secundum obiecta. Sed obiectum amoris est bonum apprehensum. Cum ergo alia sit apprehensio praesentis vitae, et alia apprehensio futurae vitae; videtur quod non maneat eadem caritas utrobique.

[36330] Iª-IIae q. 67 a. 6 arg. 3
Praeterea, eorum quae sunt unius rationis, imperfectum potest venire ad aequalitatem perfectionis, per continuum augmentum. Sed caritas viae nunquam potest pervenire ad aequalitatem caritatis patriae, quantumcumque augeatur. Ergo videtur quod caritas viae non remaneat in patria.

[36331] Iª-IIae q. 67 a. 6 s. c.
Sed contra est quod apostolus dicit, I ad Cor. XIII, caritas nunquam excidit.

[36332] Iª-IIae q. 67 a. 6 co.
Respondeo dicendum quod, sicut supra dictum est, quando imperfectio alicuius rei non est de ratione speciei ipsius, nihil prohibet idem numero quod prius fuit imperfectum, postea perfectum esse, sicut homo per augmentum perficitur, et albedo per intensionem. Caritas autem est amor; de cuius ratione non est aliqua imperfectio, potest enim esse et habiti et non habiti, et visi et non visi. Unde caritas non evacuatur per gloriae perfectionem, sed eadem numero manet.

[36333] Iª-IIae q. 67 a. 6 ad 1
Ad primum ergo dicendum quod imperfectio caritatis per accidens se habet ad ipsam, quia non est de ratione amoris imperfectio. Remoto autem eo quod est per accidens, nihilominus remanet substantia rei. Unde, evacuata imperfectione caritatis, non evacuatur ipsa caritas.

[36334] Iª-IIae q. 67 a. 6 ad 2
Ad secundum dicendum quod caritas non habet pro obiecto ipsam cognitionem, sic enim non esset eadem in via et in patria. Sed habet pro obiecto ipsam rem cognitam, quae est eadem, scilicet ipsum Deum.

[36335] Iª-IIae q. 67 a. 6 ad 3
Ad tertium dicendum quod caritas viae per augmentum non potest pervenire ad aequalitatem caritatis patriae, propter differentiam quae est ex parte causae, visio enim est quaedam causa amoris, ut dicitur in IX Ethic. Deus autem quanto perfectius cognoscitur, tanto perfectius amatur.

